\documentclass[11pt]{article}
\usepackage{fontspec}
\setmainfont{Times New Roman}
\setsansfont{Arial}
\setmonofont{Consolas}
\usepackage{amsmath,amssymb,mathtools}
\usepackage{geometry}
\usepackage{microtype}
\geometry{a4paper,margin=2.05cm}
\setlength{\parindent}{1.1em}
\setlength{\parskip}{0pt}
\newcommand{\foreign}[1]{#1}
\providecommand{\Norm}{\operatorname{N}}
\providecommand{\Gg}{\mathfrak G}
\providecommand{\Ii}{\mathfrak I}
\providecommand{\PP}{\mathfrak P}
\providecommand{\glqq}{``}
\providecommand{\grqq}{''}
\emergencystretch=220em
\allowdisplaybreaks
\sloppy
\hbadness=10000
\vbadness=10000
\hfuzz=5pt
\vfuzz=12pt
\raggedbottom
\begin{document}

\section*{41. Teorem o glavnom rodu za relativno-Galoisove čislove polja}

\begin{center}
\emph{\foreign{Math. Ann. 108 (1933), S. 411--419}}
\end{center}

V najnovějše vrěme pokazalo se, že nekomutativne metody, osoblivo teorija algeber, daju možnost formulovati i dokazati znane teoremy o relativno-cikličnyh i relativno-abelovyh čislovyh poljah za proizvoljne relativno-Galoisove polja.\footnote{Sr. moje zürihsko predavanje, \foreign{\emph{Hyperkomplexe Systeme in ihren Beziehungen zur kommutativen Algebra und Zahlentheorie}}, Verhandlungen des Internationalen Mathematiker-Kongresses Zürich 1932, Bd. I. Za teorem o razloženih algebrah sr. takodže H. Hasse, \foreign{\emph{Die Struktur der R. Brauerschen Algebrenklassengruppe}}, Math. Ann. 107 (1932).} Prěd vsim napominam normny teorem, ktory obecno možno vysloviti kako teorem o razloženih algebrah. V slědujučem pokazujem, že sootvětně i teorem o glavnom rodu možno obecno hiperkompleksno formulovati, kako teorem o vnutrnyh avtomorfizmah ili takodže o skrěženih predstavjenjah. Dokaz izhodi iz nazvanogo teorema o razloženih algebrah čistymi algebraično-aritmetičnymi razsmotrenjami, medžu tym kako v prědnem teoremu normny teorem v cikličnom slučaju, i uže pri prostom stepeni, ostaje transcendentnym jadrom.

Za bolje razuměnje formulacije najprvo prědstavljam v slědujučem neupotrěbljany teorem o glavnom rodu v minimalnom smislu,\footnote{Govorim ``v minimalnom'', bo jde o algebraičny analog teorema o glavnom rodu v proizvoljnyh poljah, medžu tym kako v ``malom'' uže imaje značenje prěhoda k \(p\)-adičnomu osnovnomu polju.} elementarno-algebraičny teorem,\footnote{Teorem nahodi se pri A. Speiser, \foreign{Zahlentheoretische Sätze aus der Gruppentheorie}, Math. Z. 5 (1919), S. 1--6; tam on uže jest vyslovjen kako teorem o skrěženih predstavjenjah.} ktory v cikličnom specialnom slučaju prěhodi v znany teorem 90 iz Hilbertovogo \foreign{Zahlbericht}, po kojem \(\Norm(a)\) jest ravno jedinici togda i samo togda, kogda \(a=b^{1-S}\). Imenno i tut teorem ja uže vyslovjam kako teorem o vnutrnyh avtomorfizmah ili o skrěženih predstavjenjah i dokazujem jego posrědstvom obćih teoremov o vnutrnyh avtomorfizmah i skrěženo-produktnom predstavjenju algeber.

Na město poslednjego v vlastnom teoremu o glavnom rodu stupa skrěženy produkt idealno-klasnoj grupy i Galoisovej grupy, ale s bolj tonkoju induciranoju idealno-klasnoju razdělboju faktornyh sistemov; tym ohvača se analog k lučno-klasnoj razdělbi teorije klasnyh polj. Kako tut obče teoremy o avtomorfizmah i predstavjenjah ješče nedostavajut, teorem o glavnom rodu možno smatrjati prvim krokom v tom směru. Zato, kako spomenuto, dajem svedenje posrědstvom teorema o razloženih algebrah. Tym teorem svodi se na sootvětny teorem za idealy vměsto idealnyh klasov; tut Brandtove teoremy o svęzi maksimalnyh porjadkov algebry\footnote{Brandtova teorija maksimalnyh porjadkov jest izložena i novo obosnovana pri H. Hasse, \foreign{Über \(p\)-adische Schiefkörper und ihre Bedeutung für die Arithmetik hyperkomplexer Zahlsysteme}, Math. Ann. 104 (1931).} možno smatrjati ekvivalentom teoremov o avtomorfizmah. One daju, v svęzi s razsmotrenjami maksimalnyh porjadkov skrěženih produktov,\footnote{Byla objavjena po jedna nota C. Chevalley, H. Hasse i E. Noether; dvě poslednje v Herbrandovom spominskom tomu.} dokaz suščstveno paralelny dokazu v minimalnom, hoti bolj komplikovany zaradi vračanja k odinokym městam. Zato prědpokladam skoro trivialny dokaz, dany od E. Artina, ktory tut prědstavlja ješče prostějši analog Speiserovogo dokaza.

\subsection*{§1. Teorem o glavnom rodu v minimalnom slučaju}

V tutom paragrafu \(k\) znači proizvoljno komutativno polje, ne nužno čislovo polje; \(K/k\) jest separabilno Galoisovo razširenje \(n\)-togo stepena, a \(\Gg\) jest sootvětnu Galoisova grupa.

Najprvo sobiraju se znane fakty o skrěženih produktah.\footnote{Teorija jest, v svęzi s mojeju lekcijeju (1929/30), izložena pri H. Hasse, \foreign{\emph{Theory of cyclic algebras}}, Trans. Amer. Math. Soc. 34 (1932), Chap. 2; sr. dalje M. Deuring o hiperkompleksnyh čislah.} Skrěženy produkt znači jednočasno vloženje \(K\) i \(\Gg\) v algebru \(A\) tako, že avtomorfizmy \(K\) stavaju se vnutrnymi. Ako \(u_{S_1},\ldots,u_{S_n}\) sut simboly k \(n\) elementam grupy, to \(A\) najprvo postavlja se kako modul linearnyh form ranga \(n\) nad \(K\):
\[
        A=u_{S_1}K+\cdots+u_{S_n}K .                         \tag{1}
\]
Posrědstvom zahtěva vnutrnogo avtomorfizma, porođenogo \(u_S\), občeje \(u_S K^*\), \(A\) stavaje se kolcom, to jest algebroju ranga \(n^2\) nad \(k\). Zahtěv vyražaje se črez
\[
        u_S^{-1}zu_S=z^S,\qquad\text{ili}\qquad
        z u_S=u_S z^S,                                      \tag{2}
\]
za vsako \(z\in K\), i črez
\[
        u_Su_T=u_{ST}a_{S,T},\qquad a_{S,T}\in K^*,           \tag{3}
\]
i črez asociativny zakon
\[
        a_{S,T}^{\,R}a_{ST,R}=a_{S,TR}a_{T,R}.                \tag{4}
\]
Ako produkt dvoh proizvoljnyh elementov definujemo črez
\[
 \Bigl(\sum_S u_S b_S\Bigr)\Bigl(\sum_T u_T c_T\Bigr)
   =\sum_{S,T}u_{ST}\,a_{S,T}\,b_S^{\,T}c_T ,
\]
to \(A\) stavaje se skrěženym produktom \(K\) s \(\Gg\) pri faktornoj sistemě \(a_{S,T}\); označenje
\[
        A=(a_{S,T},K,\Gg).
\]
Dokazuje se, že \(A\) jest prosta normalna algebra nad \(k\), to jest matrično kolco \(D_r\) \(r\)-togo stepena nad prirěđenoju dělitvenoju algebraju \(D\), i že \(K\) jest maksimalno komutativno podpolje, to jest razpadno polje. Obratno, k vsakoj takoj normalnoj dělitvenoj algebri \(D\) sut matrične kolca \(D_r\), ktore na toj način možno porođati kako skrěžene produkty.

Ako prěhodi se od \(u_S\) k \(v_S=u_Sc_S\) s \(c_S\in K^*\), nastavaju asocijovane faktorne sistemy
\[
        \bar a_{S,T}
        =a_{S,T}\frac{c_S^{\,T}c_T}{c_{ST}} .                 \tag{5}
\]
Osoblivo po (5) faktorne sistemy
\[
        \frac{c_S^{\,T}c_T}{c_{ST}}
\]
asocijovany sut k jedinici; te veličiny nazyvajut se transformacijnymi veličinami. Asocijovane faktorne sistemy sobiraju se v klasu \((a)\). Klasy \((a)\) sootvětnuju jedno-jednoznačno direktnym produktnym tvorjenjam jednoj abelovej grupy; protiv direktnomu produktu one tvore abelovu grupu, izomorfnu poelementnomu produktu klasov faktornyh sistemov. Jediničny element jest klasa razloženih algeber, to jest sistem vsih transformacijnyh veličin. Idě se o R. Brauerovu grupu klasov algeber.

Za ciklične algebry dostava se specialny slučaj: ako \(Z/k\) jest ciklično i \(S\) porođajuča substitucija, to stepenjam \(S\) sootvětnuju stepeni \(u\), i prihodi
\[
        A=Z+uZ+\cdots+u^{n-1}Z,                              \tag{1'}
\]
\[
        zu=uz^S,\qquad u^n=\alpha,                            \tag{2'--3'}
\]
pri čem
\[
        \alpha\in k^*,\qquad
        \bar\alpha=\alpha\,\Norm(c)\quad\text{za }v=uc .       \tag{4'--5'}
\]
Vsaka faktorna sistema tut sostoji iz jedinogo elementa osnovnogo polja, označenje \(A=(\alpha,Z,S)\); jedinična klasa daje se normami iz \(Z^*\), i grupa klasov algeber jest izomorfna k \(k^*/\Norm(Z^*)\).

Formulacija teorema o glavnom rodu v minimalnom slučaju upotrěblja fakt, že relacije (2)--(5) sut čisto multiplikativne i takim načinom takodže definirajut grupu razširenja \(\Gg^*\) grupy \(\Gg\), sostavljenu iz elementov \(n\) kompleksov \(u_S K^*\):
\[
        \Gg^*=\{\,u_{S_1}K^*,\ldots,u_{S_n}K^*\,\}.            \tag{6}
\]
\(\Gg^*\) imaje \(K^*\) kako abelov normalny podgrup, medžu tym kako \(\Gg^*/K^*\simeq\Gg\). \(\Gg^*\) sostoji iz vsih regularnyh elementov \(g^*\), ktore transformujut \(K\) v sebe, to jest \(g^{*-1}Kg^*=K\). Kolcove avtomorfizmy \(K\) sut zato porođane vsimi i samo elementami \(\Gg^*\). Bo ako \(v=uw\) porođaje vnutrny avtomorfizm \(K\), togda \(w\) komutuje so vsimi elementami \(K\), znači leži v \(K^*\), poněže \(K\) jest maksimalno komutativno podkolco.

\paragraph{Teorem o glavnom rodu v minimalnom slučaju.}
Prva formulacija: vsak grupovy avtomorfizm \(\Gg^*\), ktory jest prodolženjem identičnogo avtomorfizma \(K^*\), jest vnutrny i porođaje se elementom iz \(K^*\). Drugymi slovami: ako transformacijne veličiny
\[
        \frac{c_S^{\,T}c_T}{c_{ST}}
\]
vse imajut vrědnost jedinica, togda \(c_S\) sut simbolične \((1-S)\)-te stepeni,
\[
        c_S=b^{1-S}=\frac{b}{b^S}\qquad(S\in\Gg).
\]
Treta formulacija: grupa \(\Gg\) imaje samo jedinu skrěženu klasu predstavjenja prvogo stepena v \(K^*\), naležeču k faktornoj sistemě jedinica.

Predstavjenje \(u_S\mapsto C_S\) nazyva se skrěženo s faktornoju sistemoju \(a_{S,T}\), ako
\[
        C_S^{\,T}C_T=C_{ST}a_{S,T}
\]
važy. Dvě predstavjenja naležet k toj samoj klasi, ako
\[
        C_S=B^{-S}D_SB .
\]

Dokazujem prvu formulaciju i potom pokazujem jej soglasje s drugoju i tretoju. Dokaz opiraje se na tom, že vsak grupovy avtomorfizm označenogo roda možno prodolžiti do kolcovogo avtomorfizma \(A\), ktory, kako znano, jest vnutrnym avtomorfizmom. Nehaj pri tom avtomorfizmu \(v_S\) sut obrazy \(u_S\). Togda \(\sum v_SK\) znovu tvori skrěženy produkt \(\Gg\) s \(K\) k toj samoj faktornoj sistemě, bo po predpostavjenju vse relacije (2)--(4) ostavajut. Prirěđenje
\[
        \sum u_Sb_S\longmapsto \sum v_Sb_S
\]
predstavja kolcovy avtomorfizm, pri kojem elementi \(K\) ostavajut fiksne. Zato on porođaje se elementom iz \(K^*\).

Prěhod k drugoj formulaciji opiraje se na tom, že po (2) \(v_S\) moraju iměti formu \(v_S=u_Sc_S\); kako i faktorne sistemy ostavajut, po (5) sootvětne transformacijne veličiny moraju stavať se jediniceju. Obratno, (2)--(5) pokazyvajut, že substitucija \(v_S=u_Sc_S\) porođaje avtomorfizm označenogo roda. Po prvoj formulaciji zato dostava se
\[
        v_S=b^{-1}u_Sb=u_S\,\frac{b}{b^S},
        \qquad c_S=b^{1-S}.
\]
V cikličnom slučaju predpostavjenje osoblivo znači \(\Norm(c)=1\), i \(c=b^{1-S}\) jest znany teorem. Treta formulacija jest neposrědnje ravnoznačna drugoj, bo predpostavjenje govori, že \(u_S\mapsto c_S\) tvori skrěženo predstavjenje prvogo stepena k faktornoj sistemě jedinica; \(c_S=b^{1-S}\) znači ekvivalenciju k jediničnomu predstavjenju.

\subsection*{§2. Teorem o glavnom rodu}

\paragraph{Prědvaritelne zamětky o klasnoj razdělbi.}
V teoremu o glavnom rodu teorije klasnyh polj za ciklične relativne polja javjajut se, kako znano, dvě različne klasne razdělby: absolutne idealne klasy v gornjem polju, lučne klasy v dolnjem polju. Tomu v obćem slučaju sootvětuje, že klasna razdělba idealov gornjego polja induciraje tonšu idealno-klasnu razdělbu faktornyh sistemov.

Nehaj najprvo \(\Ii\) znači grupu vsih idealov iz \(K\). Kako \(\Gg\) porođaje avtomorfizmy v \(\Ii\), razširenje s \(\Gg\) definirano jest relacijami (2)--(5) i sostoji iz \(n\) kompleksov \(\{\,\cdots u_S\Ii\cdots\,\}\). Ako v \(\Ii\) prěhodi se črez ustanovjenje ekvivalencije k grupě absolutnyh idealnyh klasov, to i teper \(n\) kompleksov \(u_S\bar\Ii\) sut definirane. Bo \(\Gg\) porođaje takodže avtomorfizmy klasnoj grupy, poněže glavna klasa transformuje se v sebe. Tym ekvivalencijny odnos od \(\Ii\) k \(\bar\Ii\) prěnosi se od \(\{\,\cdots u_S\Ii\cdots\,\}\) na \(\{\,\cdots u_S\bar\Ii\cdots\,\}\). Osoblivo \(u_S\) stavaje ekvivalentno vsim produktam \(u_S(c_S)\), kde \((c_S)\) proběgaje glavne idealy.

Ako teper zahtěvamo, že produktne relacije (3) sut jednoznačne v smislu toj ekvivalencije, to znači, že vse transformacijne veličiny iz glavnyh idealov stavaju jednoznačno ekvivalentne. Drugymi slovami:

V sistemě \(n\) kompleksov \(\{\,\cdots u_S\bar\Ii\cdots\,\}\), kde \(\bar\Ii\) jest grupa absolutnyh idealnyh klasov iz \(K\), a \(H\) osoblivo znači glavnu klasu, produktne relacije (3) za \(u_SH\) sut togda i samo togda jednoznačno definirane, kogda v induciranoj idealno-klasnoj razdělbi faktornyh sistemov glavna klasa soderži vse transformacijne veličiny iz glavnyh idealov.

\paragraph{Definicija.}
Inducirana idealno-klasna razdělba faktornyh sistemov ustanavlja se tak, že glavna klasa sostoji iz vsih glavnyh idealov \((a_{S,T})\), za ktore s najmanje jednym sistemom bazisnyh elementov \(a_{S,T}\) algebra
\[
        (a_{S,T},K,\Gg)
\]
razloži se na vsih razvětvjenyh městah \(K\).

Ti glavne idealy tvore grupu po spojivanju faktornyh sistemov; ona obhvata transformacijne veličiny
\[
        \frac{(c_S)^{\,T}(c_T)}{(c_{ST})},
\]
bo faktorne sistemy \(c_S^{\,T}c_T/c_{ST}\) porođajut povsudu razložene algebry.

\paragraph{Formulacija teorema o glavnom rodu.}
Teorem, sootvětně minimalnomu teoremu, vyslovja se v trěh ravnoznačnyh formah.

Prva formulacija: ako pri osnovanoj induciranoj idealno-klasnoj razdělbi faktornyh sistemov substitucija
\[
        v_S=u_S c_S
\]
porođaje avtomorfizm \(n\) kompleksov \(\{\,\cdots u_S\bar\Ii\cdots\,\}\), to jest ako za \(v_S\) v smislu toj klasnoj razdělby važet te same relacije (3), togda toj avtomorfizm jest vnutrny i porođaje se idealnoju klasoju \(b\).

Druga formulacija: ako transformacijne veličiny, stvorjene iz idealnyh klasov \(c_S\),
\[
        \frac{c_S^{\,T}c_T}{c_{ST}}
\]
vse naležet k glavnoj klasi induciranoj klasnoj razdělby faktornyh sistemov, togda vektory \(\{\,\cdots c_S\cdots\,\}\) iz idealnyh klasov tvore glavny rod, a klasy \(c_S\) sut simbolične \((1-S)\)-te stepeni; to jest jest idealna klasa \(b\) taka, že
\[
        c_S=b^{1-S}=\frac{b}{b^S}\qquad(S\in\Gg).
\]

Treta formulacija: pri osnovanoj induciranoj idealno-klasnoj razdělbi za faktorne sistemy grupa \(\Gg\) imaje samo jedinu skrěženu klasu predstavjenja prvogo stepena v idealno-klasnoj grupě, naležeču k jediničnoj klasi faktornyh sistemov.

Že te formulacije sut ravnoznačne, slěduje kako v §1. Prěhod od prvoj k drugoj formulaciji tut jest čak prostějši, bo avtomorfizm uže jest predpostavjen kako porođeny črez \(v_S=u_Sc_S\). K tretoj formulaciji treba samo zamětiti, že predstavjenja v idealno-klasnoj grupě sut definirane samo multiplikativno; izjava o predstavjenjah prvogo stepena tym ne dobivaje ograničenja.

\paragraph{Dokaz teorema o glavnom rodu.}
Dokaz drugoj formulaciji opiraje se na dvě pomočne teoremy.

\paragraph{Pomočny teorem 1.}
Ako transformacijne veličiny, stvorjene iz idealov \(c_S\),
\[
        \frac{c_S^{\,T}c_T}{c_{ST}}
\]
daju jediničny ideal, togda \(c_S\) sut simbolične \((1-S)\)-te stepeni:
\[
        c_S=b^{1-S}\qquad(S\in\Gg).
\]
Ravnoznačno jest to s prvoju formulacijeju, ako tam predpostavlja se, že avtomorfizm porođaje se črez \(v_S=u_Sc_S\). Dokaz: \(b=\sum c_S\) daje požadovano, pri čem suma jest modulna suma, to jest največši občy dělitelj. Bo
\[
        b=\sum_S c_S,\qquad
        b^T=\sum_S c_S^{\,T},\qquad
        b^T c_T=\sum_S c_S^{\,T}c_T=\sum_S c_{ST}=b,
\]
tako \(c_T=b^{1-T}\).

\paragraph{Pomočny teorem 2.}
Ako algebra
\[
        A=(a_{S,T},K,\Gg)
\]
razloži se na vsih konečnyh i beskonečnyh razvětvjenyh městah \(K\), i glavne idealy \((a_{S,T})\) sut transformacijne veličiny idealov,
\[
        (a_{S,T})=\frac{c_S^{\,T}c_T}{c_{ST}},
\]
togda \(A\) razloži se povsudu, to jest bezuslovno. Togda \((a_{S,T})\) stavaju transformacijne veličiny glavnyh idealov,
\[
        (a_{S,T})=\frac{(d_S)^{\,T}(d_T)}{(d_{ST})}.
\]
Obratno, vsaka povsudu razložena algebra zadovolja tomu uslovju.

Dokaz opiraje se na znanom povedanju \(A\) pri \(p\)-adičnom razširenju. Nehaj \(k_p\) znači \(p\)-adično razširenje \(k\), sootvětně \(K_\PP\) \(\PP\)-adično razširenje \(K\), kde \(\PP\) jest prosty ideal nad \(p\); \(A_p\) nehaj bude razširenje \(A\), nastalo prěhodom od \(k\) k \(k_p\). Dalje nehaj \(\mathfrak Z\) bude grupa razloženja pri \(\PP\), a \(a_{\mathfrak Z}\) sootvětny izrěz faktornoj sistemy. Togda važy
\[
        A_p\sim (a_{\mathfrak Z},K_\PP/k_p,\mathfrak Z).
\]
Ako \(p\) jest osoblivo nerazvětvjeny, znači \(\mathfrak Z\) ciklična, togda idě o vyznačeno predstavjenje \(A_p\) posrědstvom nerazvětvjenogo inercijnogo polja \(K_\PP/k_p\).

Treba pokazati, že pod predpostavjenjem pomočnogo teorema 2 \(A_p\) takodže razloži se v tom nerazvětvjenom slučaju. Za beskonečne města tut \(K_\PP=k_p\), zato \(A_p\) razloži se. Za konečno \(p\) najprvo slěduje, že faktorna sistema \(a_{\mathfrak Z}\) asocijovana jest s faktornoju sistemoju, sostavjenoju iz jedinic. Pri prěhodu k normovanomu cikličnomu predstavjenju v \(A_p=(\alpha,K_\PP/k_p,R)\), s \(R\) kako porođajučim elementom \(\mathfrak Z\), faktorna sistema daje se relacijami
\[
        u_{\PP_i}u_{\PP_j}=u_{\PP_{i+j}}\varepsilon_{\PP_i,\PP_j}
\]
Stavivši \(u_R=u\), dostava se
\[
        u^i=u_R^i\,\varepsilon_R\varepsilon_{R^2}\cdots \varepsilon_{R^{i-1}},
\]
tako osoblivo \(u_E=\varepsilon_E\), kogda \(E\) jest porjadok \(\mathfrak Z\). Poněže v nerazvětvjenom razširenju \(K_\PP/k_p\) vse jedinice sut normy, slěduje, že \(A_p\) razloži se. Zato \(A\) razloži se povsudu, i tedy po teoremu o razloženih algebrah bezuslovno.

Drugymi slovami: iz predpostavjenja slěduje, že \(a_{S,T}\) sut transformacijne veličiny iz elementov,
\[
        a_{S,T}=\frac{d_S^{\,T}d_T}{d_{ST}},
        \qquad d_S\in K^*.
\]
Oslabjeno to znači: \((a_{S,T})\) stavaju transformacijne veličiny glavnyh idealov. Ako obratno \(A\) jest povsudu razložena, togda i glavne idealy \((a_{S,T})\) stavaju transformacijne veličiny idealov na vsakom městu, i \(A\) razloži se faktično po drugoj formulaciji predpostavjenja teorema o razloženih algebrah.

Iz dvuh pomočnyh teoremov neposrědnje slěduje dokaz teorema o glavnom rodu. Nehaj \(c_S\) znači nekaky ideal iz klasy \(c_S\). Predpostavjenje govori, po definiciji induciranoj klasnoj razdělby, že transformacijne veličiny \(c_S\) stavaju glavne idealy \((a_{S,T})\), ktore zadovoljajut predpostavjenja pomočnogo teorema 2. Zato sut glavne idealy \((d_S)\), tako že
\[
        \frac{c_S^{\,T}c_T}{c_{ST}}
\]
črez \((d_S)\) stavaje jedinicoju. Idealy \(c_S/(d_S)\), ktore ležet v toj samoj klasi \(c_S\), takim načinom zadovoljajut predpostavjenje pomočnogo teorema 1. Zato jest idealna klasa \(b\) s
\[
        c_S=b^{1-S}.
\]
To znači, že klasy \(c_S\) sut simbolične \((1-S)\)-te stepeni klasy \(b\).

V cikličnom slučaju teorem o glavnom rodu prěhodi v znany teorem, kogda glavna klasa induciranoj klasnoj razdělby pri osnovanoj normaciji prěhodi v grupu tyh glavnyh idealov \((\alpha)\), za ktore \(\alpha\) na vsih razvětvjenyh městah stavaje normnym ostatkom.

\begin{center}
(Dostavljeno 27. 10. 1932.)
\end{center}

\end{document}
